\chapter{What is a light-field microscope}
\label{ch:light-field-microscope}

In this chapter, a physical model describes how the fluorescent specimen produces a sensor frame and how sub-aperture images (views) are extracted from that frame.
The computational target is a single two-dimensional image formed by aligning and combining those views, as described in Chapter~\ref{ch:image-reconstruction}; the camera directly measures the sensor frame.
The model will be used to describe the formed images, the system's characteristics, and its limitations.
This chapter also describes the optical degradation mechanisms and their effect on scene recovery.
The model is continuous in space, while sensor and simulation outputs are discrete and therefore require sampling steps.
The specifics describing the physical traits of the optical system are discussed in \cref{ch:experimental-setup}.

In microlens-based LFM, a microlens array placed at an image plane allows the camera to record a two-dimensional array of local angular measurements in a single exposure. The resulting four-dimensional light field can be computationally resampled to synthesize viewpoints, refocus after acquisition, or, in suitable systems, reconstruct a three-dimensional volume \cite{levoy2006lightFieldMicroscopy,levoy2009recordingControlling4d,broxton2013waveOpticsLfm}. LFM does not eliminate the imaging trade-offs; it changes where they occur. The finite sensor and optical aperture impose a spatial-angular sampling trade-off: allocating more samples to angular diversity leaves fewer spatial samples per view, while the resulting sampling pattern affects the resolution and quality of computational reconstruction.

Fourier light-field microscopy (FLFM) instead places the microlens array in a Fourier, or pupil-conjugate, plane of the imaging system. This arrangement records the spatial and angular components of the light field with a spatially invariant sampling scheme and permits a unified three-dimensional PSF under the system model described by Guo et al. \cite{guo2019fourierLightField}.

The physical system is a Fourier light-field fluorescence microscope, in which the specimen is excited by a $488\,\mathrm{nm}$ laser and the emitted fluorescent light is collected through an objective.
The specimen's emission is distinct from the detected intensity, the optical quantity used in this model: contributions from mutually incoherent fluorescent emitters add as intensities, as expressed in the FLFM image-formation model of Guo et al. \cite{guo2019fourierLightField}.
Camera values are digital counts resulting from integration over each pixel area and exposure time, followed by detector conversion.
Spectral integration, exposure duration, gain, and offsets must be accounted for in that conversion; the detected emission band and camera calibration remain to be specified in \cref{ch:experimental-setup}.

To orient the image-formation model, $(x_o,y_o)$ denote continuous spatial coordinates on the nominal specimen plane, with their origin on the optical axis; $z=0$ is this plane and positive $z$ points toward the collection objective.
The coordinates $(u,v)$ denote physical positions on the plane conjugate to the objective's back pupil, with their origin at the pupil center and axes chosen along the microlens grid.
Together, the spatial and pupil coordinates provide geometric bookkeeping for position and collection direction, consistent with the spatial-directional interpretation of microscope light fields \cite{levoy2009recordingControlling4d}.
They are continuous length coordinates, distinct from the discrete view and pixel indices introduced next.

The image formed on the camera sensor is built from a matrix of sub-aperture images, also called elemental view (or views, for short), each corresponding to a specific microlens in the microlens array.
Let $Y(r_s,c_s)$ denote the raw sensor frame in sensor row and column indices, separately from the extracted view $V_{ij}(y,x)$.
Following Chapter~\ref{ch:image-reconstruction}, we use one-based view indices $i,j\in\{1,\ldots,N\}$ and within-view row and column coordinates $(y,x)\in\Omega=\{1,\ldots,H\}\times\{1,\ldots,W\}$.
The extracted light-field stack has shape $(N,N,H,W)$, where $N$ counts retained views along each view axis, rather than all manufactured lenslets, and $H$ and $W$ are each view's height and width in pixels.
Writing $V=\mathcal E(Y)$, $\mathcal E$ denotes extraction using the selected sensor windows; this operation follows camera integration.
The first view index advances horizontally on the sensor and the second vertically; the one-based indices map to zero-based array indices by subtracting one.
The extraction geometry and array-axis mapping will be documented in \cref{sec:sample-acquisition}.

The exact path of light from the specimen to the different sensors depends on the specific capture discussed.
There are 3 capture modes: (1) the light field path, (2) the traditional imaging path, and (3) the validation path.
This chapter focuses on the light field path, which is the main path of interest for this work, while the others are used as comparison tools.
The light field path can be broken down into 3 sequential steps: first, the light emitted is collected by the objective and the Fourier plan is formed (or the back focal plane of the objective); second, the light is magnified by a simple 4F system onto the microlens array which forms the light field views; and third, the light field is de-magnified by another simple 4F system to fit inside the camera sensor.
It is important to note that the camera is a discrete, intensity-based sensor, so the actual light-field data measured is a discrete sampling of the intensity of the light field (which is a complex continuous function in space).

Beyond the image formation described above, the specimen itself degrades the measurement. In biological tissue, absorption and scattering reduce penetration and distort the detected light, with the severity depending on the sample and wavelength \cite{jacques2013opticalPropertiesTissues}. This limits the depth and contrast achievable by purely optical systems, especially when the sample is thick, heterogeneous, or strongly scattering.



